\section{Design}
\label{sec:design}

extC keeps the part of a C-like language that makes memory behaviour predictable:
objects never move, layouts are deterministic (there is no object header, and no word of a value
is ever reserved for the runtime), and there are no write barriers on data stores.
The design question we answer here is how reclamation fits into that shape without giving up
bounded memory or bounded pause.

Our answer has three parts.
(1) \emph{Storage is stratified by lexical scope.} Every allocation point has a statically known
level, and the level selects one of five substrates (\S\
ef{sec:substrates}), each with a different
reclamation rule.
(2) \emph{Reclamation is a compile-time decision.} A unit is analysed and placed in one of three
tiers (\S\
ef{sec:tiers}): no runtime at all, an event tier where reference counts free objects on
the spot, or a sweep tier where each region is marked and swept under explicit budgets.
(3) \emph{The runtime obeys a two-dimensional contract.} Memory is bounded by a wall that the
collector's trigger line respects, and each synchronous step of the collector is bounded by a
budget that is independent of the heap size (\S\
ef{sec:contract}).

\begin{table}[t]
\centering
\small
\begin{tabular}{@{}p{0.15\linewidth}p{0.19\linewidth}p{0.26\linewidth}p{0.28\linewidth}@{}}
\toprule
Substrate & Lifetime owner & Allocation & Reclamation \\
\midrule
Frame arena & lexical scope of a block & bump pointer in a chunked arena & whole arena rewound on scope exit \\
RC region & frame level (a function or block) & slot in a slab, size class by request & budgeted mark--sweep inside the region \\
Pool & the \texttt{place} the container was created in & slab chain rooted at that place & chain rewound with the place \\
Task arena & coroutine frame or parallel worker & bump, thread-local & arena discarded when the task ends \\
Plate & explicit \texttt{plate} value, OS pages & \texttt{mmap} of page-rounded bytes & \texttt{mprotect} / \texttt{munmap} \\
\bottomrule
\end{tabular}
\caption{The five storage substrates. Only the third row is managed by the reclamation runtime;
the others are released wholesale, so their cost is paid once per scope or per task.}
\label{tab:substrates}
\end{table}

\subsection{Storage substrates}
\label{sec:substrates}

\tab\ref{tab:substrates} lists the substrates.
Which one a \texttt{new} lands on is decided by the level of its frame, not by its type, so the
same struct can be arena-allocated in one place and region-allocated in another.

\paragraph{Frame arenas.}
A block that is not a reclamation region gets a bump arena: a chunk list whose allocation is one
pointer increment, and whose exit restores the pointer to the block entry.
The zero cost of this row is what the whole design is trying to preserve, and \S\
ef{sec:eval}
shows that 71.7\% of the compiling programs in our corpus still get it for free (\S\
ef{sec:census}).

\paragraph{Reclamation regions.}
A region is a set of slabs owned by one frame level.
Each object occupies exactly one slot; the slot stride is the request rounded up to a multiple of
eight, so an 80-byte object occupies 80 bytes rather than the next power of two.
Per-slot metadata is kept in three bitmaps (occupied, marked, cross-region registration) plus a
16-bit type number and a 16-bit in-edge count, so the payload is byte-exact and the layout of a
value is never perturbed by the runtime (\S\
ef{sec:impl-geometry}).

\paragraph{Pools and plates.}
Containers whose element type is known allocate from a slab chain attached to the \emph{place}
where they were created, so a container that dies with a block releases its storage with that
block instead of being swept.
Plates are the escape hatch for data that must outlive the process's own structure: page-aligned
\texttt{mmap}, optionally protected, released explicitly.

\paragraph{Task arenas.}
A coroutine frame and a parallel worker each get a thread-local arena that is discarded when the
task ends (\texttt{extc\_par\_close} for workers, frame teardown for coroutines).
Freshly spawned tasks that write into their owner's view get references pinned to the owner's
region instead, so a task can never leave a dangling edge behind when it is torn down.

\subsection{Placement and escape}
\label{sec:placement}

Levels are assigned by the escape analysis: parameters are at level $0$, a function body at
$1$, and each nested block adds one.
An allocation point's region is the level where its storage must live:
\begin{equation}
\ell(pi) = \max(0, k - \mathit{escapeDepth}(pi)),
\label{eq:placement}
\end{equation}
where $k$ is the lexical depth of the allocation point and $escapeDepth$ is the number of
levels the value escapes.
Three rules keep the stratified store sound:

\begin{itemize}
\item \emph{D1 (no downward edge).} Storing a reference into a slot at depth $e$ requires the
  pointed-to storage to be at depth $d \le e$. Storing a younger object into an older slot is
  allowed; the reverse would let a scope outlive its referent.
\item \emph{D2 (promotion is static).} If a value may escape to level $d$, its allocation point
  is placed at level $d$ at compile time: the storage is born there, nothing is copied or moved at
  run time, and the object's address is fixed from birth. The only run-time cost at an escape site
  is the one pointer store that puts the reference into an outer slot.
  Whether a value lives in a long-lived region is therefore known statically, and
  \S\ref{sec:formal} never has to reason about a move.
\item \emph{D3 (views).} A view may not outlive the storage it views: $d \le c$ for view
  $c$ of container $d$.
\end{itemize}

Promotion is where unbounded memory can appear: the storage now lives until the outer level ends,
although only one of the values written into it may still be reachable. That is the shape the
reclamation tier exists for, and the shape \S\
ef{sec:formal} reasons about.

\subsection{Choosing a tier}
\label{sec:tiers}

After placement, the unit is analysed as a whole and assigned a tier.
The assignment is a pure compile-time decision; two user-visible switches
(\texttt{EXTC\_RECLAIM}, \texttt{EXTC\_RC\_SWEEPFREE}) only force or forbid a tier, and they exist
for the experiments in \S\
ef{sec:eval}.

\paragraph{Tier 0: no runtime.}
If no allocation point is governed by reclamation, the emitted program contains no reclamation call
and the runtime text is not even emitted (a 434-line \texttt{hello} world has zero
\texttt{extc\_rc\_} tokens). This is the default outcome for programs that never promote.

\paragraph{Idle units.}
A unit is \emph{idle} when every reference field is written at most once per object, no reference
local or parameter is reassigned in a way that drops a reference, and no operation drops a
reference (\S\
ef{sec:formal-idle} formalises this).
An idle unit needs no reclamation work at all: objects only disappear when their whole region is
released. The judgement is conservative and whole-unit, so one function that fails it keeps the
runtime for the unit.

\paragraph{The event tier (T1).}
The event tier is taken when the unit is \emph{acyclic}: every reference write either stores a
freshly allocated value, \texttt{null}, or a value from a strictly older level, and every slot that
can drop a reference is a counted slot.
Under those conditions no cycle can form and every object that becomes unreachable reaches count
zero, so the runtime keeps reference counting and does not emit sweeping.

\paragraph{The sweep tier (T2).}
Otherwise the unit gets cyclic-capable regions: counting is off by default (the
\emph{no-count} configuration) and unreachable storage is recovered by budgeted mark--sweep over one
region at a time.
Counting is not needed for safety in this tier -- the marker finds its seeds from the static root
table and the cross-region registration bits, not from counts -- so switching it off only removes a
fast path, never a guarantee.

\begin{table*}[t]
\centering
\small
\resizebox{\ifdim\width>\linewidth\linewidth\else\width\fi}{!}{\begin{tabular}{@{}llccc@{}}
\toprule
Tier & When & Counting & Sweeping & Contract \\
\midrule
T0 & user forces \texttt{EXTC\_RECLAIM=0} & no & no & memory not bounded (control group) \\
none & $\mathit{idle}(U)$ & no & no & region wholesale release only \\
T1 & $\mathit{acyclic}(U)$ & yes & no & $O(1)$ free, memory $\le L + G + O(B)$ \\
T2 & otherwise & no (default) & yes & budgeted mark--sweep \\
\bottomrule
\end{tabular}}
\caption{Tier assignment. The last two rows are the ones that emit a runtime; T0 and the idle case
emit none. \S\
ef{sec:formal} gives the soundness statement for each row.}
\label{tab:tiers}
\end{table*}

\subsection{The two-dimensional contract}
\label{sec:contract}

Let $L$ be the live set of a region (bytes reachable from a complete root table), $M$ the
bytes the region has ever admitted, and $M_max$ the wall.
Two properties are promised for every program that runs in T1 or T2:

\begin{description}
\item[Bounded memory.] $M \le L + G + O(B)$ where $G$ is the trigger floor and $B$
  the size of one slab, and the bound does not depend on the number of loop iterations, i.e.~it is a
  function of the live set, not of the total allocation volume.
\item[Bounded pause.] Every synchronous step the collector takes on the mutator thread is bounded by
  a constant $P_{\max}$ fixed at compile time, independent of the heap size and of the size of
  the object being scanned. Pause here means the sum of all collector time on that thread, not a
  stop-the-world window: there is no separate collector thread to hide work in.
\end{description}

The contract is maintained by four release layers, each of which removes a different class of
garbage:
\begin{enumerate}
\item \emph{Scope exit} rewinds the arena (zero cost).
\item \emph{Static placement} puts the storage of an escaping value directly at the level it
  escapes to; at run time an escape costs one pointer store, never a copy.
\item \emph{Count zero} frees exactly one object at the moment its last holder drops it (T1, and
  optionally T2).
\item \emph{Region sweep} recovers what the first three cannot: cycles, and storage whose counts were
  saturated deliberately.
\end{enumerate}

Counting is a knob, not a guarantee: turning it off changes which layer reclaims a given object, not
whether the program is safe, and the sweep tier therefore stays correct with counting disabled.

\subsection{Deliberate non-goals}
\label{sec:non-goals}

We do not move objects, so no compaction and no copying collection: fragmentation is bounded by the
size-class layout, not eliminated.
We do not scan stacks: the root table is static, emitted at each trigger point, which is why the
collector needs no shadow stack and why a suspended coroutine registers its frame fields instead.
We do not run the collector on another thread, so \emph{every} collector instruction is mutator time;
this is the honest reading of our pause numbers and the reason \S\
ef{sec:eval} reports both the
armed-window total and the per-step total.
We do not attempt to bound the pause of a program that allocates an object larger than the wall in a
single request: the admission rule (\S\
ef{sec:impl-admission}) reserves room for one maximum-sized
request, so such a program is rejected by the bound rather than silently breaking it.

\begin{table*}[t]
\centering
\small
\resizebox{\ifdim\width>\linewidth\linewidth\else\width\fi}{!}{\begin{tabular}{@{}lll@{}}
\toprule
Property & Status & Mechanism or reason \\
\midrule
peak memory & bounded by $L+G+O(B)$ & trigger line $\ge 2L$, wall $M_{\max}$ \\
per-step collector time & bounded & batch cap 256, segment deadline 40\,\us \\
region destruction & $O(1)$ & closing queue drained at allocation points \\
one request larger than wall & reserved & admission: $\mathit{thr} + R_{\max} \le M_{\max}$ \\
long-lived garbage in a region & bounded, not zero & sweep is triggered by pressure, not by time \\
count saturation & not a leak & saturated slots stay countable-by-sweep \\
pause of a single huge scan & bounded & object scan chunked into 512-word segments \\
\bottomrule
\end{tabular}}
\caption{What the current design does and does not bound. Each bounded row is backed by a theorem
in \S\
ef{sec:formal} or by a measured gate in \S\
ef{sec:eval}.}
\label{tab:bounded}
\end{table*}